\chapter{The continuous-flow multiplicative ergodic theorem}
\label{chap:continuous}

\section{Overview}

The discrete multiplicative ergodic theorem governs a single measure-preserving
map \( T : X \to X \) and the iterates of a matrix cocycle generated one step at a
time. This chapter lifts that theorem to \emph{continuous time}: the map \( T \) is
replaced by a one-parameter measure-preserving flow \( \varphi : \R \to X \to X \),
and the iterated cocycle by a continuous-time linear cocycle
\( A : \R \to X \to \mathrm{Mat}_d(\R) \). The conclusion is a finite Lyapunov
spectrum \( \lambda_1 > \cdots > \lambda_k \), a measurable filtration that is
\emph{flow-equivariant} at every real time, and the exact continuous-parameter
growth \( t^{-1} \log\norm{A(t,x)v} \to \lambda_i \) as \( t \to \infty \) over \( \R \).

The strategy is a reduction, not a redevelopment of the ergodic machinery for \( \R \).
We set \( T := \varphi(1) \) and read the discrete cocycle off the sampled flow cocycle;
the proved discrete theorem \ref{oseledets_filtration} delivers the integer-time
conclusion, and two analytic devices lift it to the continuous parameter: a
\emph{between-times sandwich} that controls growth on each interval \( [n, n+1) \), and a
\emph{shift-invariance of the growth} \( \limsup \) that promotes the discrete (time-one)
equivariance to equivariance at every real time. No continuous-time Kingman theorem is
needed; the integer clock appears only as a technical reduction device. Throughout,
\( X \) carries no topology: the flow is a measurable measure-preserving action of
\( (\R, +) \), and the cocycle is measurable in the state at each fixed time.

\section{The continuous-time data}

\begin{definition}[Measure-preserving one-parameter flow]\label{MeasurePreservingFlow}
  \lean{Oseledets.MeasurePreservingFlow}\leanok
  A \emph{measure-preserving one-parameter flow} on a measurable space \( X \) for a measure
  \( \mu \) is a family \( \varphi : \R \to X \to X \) together with the data
  \( \varphi(0) = \mathrm{id} \), \( \varphi(s+t) = \varphi(s) \circ \varphi(t) \) for all
  \( s, t \in \R \), and a proof that each time-\( t \) map \( \varphi(t) \) preserves \( \mu \).
  No topology on \( X \) is assumed; in particular no joint continuity in \( (t, x) \) is
  required.
\end{definition}

\begin{lemma}[Integer times are iterates of the time-one map]\label{natCast_eq_iterate}
  \lean{Oseledets.MeasurePreservingFlow.natCast_eq_iterate}\leanok
  For a measure-preserving flow \( \varphi \) and \( n \in \N \), the integer-time map of the
  flow is the \( n \)-fold iterate of its time-one map: \( \varphi(n) = (\varphi(1))^{[n]} \).
\end{lemma}
\begin{proof}\leanok
  \uses{MeasurePreservingFlow}
  Induction on \( n \). The base case \( \varphi(0) = \mathrm{id} = (\varphi(1))^{[0]} \) is the
  time-zero law. For the step, write \( n+1 = n + 1 \) as reals and apply additivity
  \( \varphi(n+1) = \varphi(n) \circ \varphi(1) \), then the inductive hypothesis and
  \( (\varphi(1))^{[n+1]} = (\varphi(1))^{[n]} \circ \varphi(1) \).
\end{proof}

\begin{definition}[Continuous-time linear cocycle over a flow]\label{FlowCocycle}
  \lean{Oseledets.FlowCocycle}\leanok
  A \emph{continuous-time linear cocycle} over a measure-preserving flow \( \varphi \), valued in
  invertible \( d \times d \) real matrices, is a family
  \( A : \R \to X \to \mathrm{Mat}_d(\R) \) with \( A(0,x) = 1 \), the cocycle identity (newest
  factor on the left) \( A(t+s,x) = A(t, \varphi(s)x)\, A(s,x) \), a proof that
  \( \det A(t,x) \ne 0 \) for all \( t, x \), and measurability of each time-\( t \) map
  \( A(t,\cdot) \).
\end{definition}

\begin{proposition}[Reduction identity at integer times]\label{toCocycle_eq}
  \lean{Oseledets.FlowCocycle.toCocycle_eq}\leanok
  For a flow cocycle \( A \) over \( \varphi \), every \( n \in \N \) and every \( x \),
  \[ A(n, x) = \mathrm{cocycle}\,(A(1,\cdot))\,(\varphi(1))\, n\, x, \]
  i.e.\ at integer times the continuous-time cocycle equals the discrete iterated cocycle
  generated by its time-one map \( A(1,\cdot) \) over the time-one dynamics \( \varphi(1) \).
\end{proposition}
\begin{proof}\leanok
  \uses{FlowCocycle,natCast_eq_iterate}
  Induction on \( n \). At \( n = 0 \) both sides are the identity. For the step, split
  \( A((n+1),x) = A(n, \varphi(1)x)\, A(1,x) \) by the cocycle identity at \( t = n \), \( s = 1 \),
  match the recursion \( \mathrm{cocycle}\,(n+1)\,x = \mathrm{cocycle}\, n\, (\varphi(1)x) \cdot A(1,x) \),
  and apply the inductive hypothesis at the point \( \varphi(1)x \).
\end{proof}

\section{Reduction to the discrete theorem}

The discrete theorem requires the time-one generator to have integrable positive log-norm,
both forward and inverse. These follow by evaluating the uniform dominating hypotheses at
\( s = 1 \).

\begin{lemma}[Integrability of the time-one log-norm]\label{integrableLogNorm_timeOne}
  \lean{Oseledets.integrableLogNorm_timeOne}\leanok
  If \( g \in L^1(\mu) \) dominates \( \log^{+}\norm{A(s,x)} \) for all \( s \in [0,1] \) and all
  \( x \), then the time-one map \( A(1,\cdot) \) has integrable positive log-norm.
\end{lemma}
\begin{proof}\leanok
  \uses{FlowCocycle}
  The map \( x \mapsto \log^{+}\norm{A(1,x)} \) is measurable (composition of the measurable
  \( \log^{+} \), the operator norm, and the measurable time-one cocycle map) and is dominated
  pointwise by \( g \), taking \( s = 1 \in [0,1] \) in the hypothesis. Dominated by an integrable
  function, it is integrable.
\end{proof}

\begin{lemma}[Integrability of the inverse time-one log-norm]\label{integrableLogNorm_timeOne_inv}
  \lean{Oseledets.integrableLogNorm_timeOne_inv}\leanok
  If \( g' \in L^1(\mu) \) dominates \( \log^{+}\norm{(A(s,x))^{-1}} \) for all \( s \in [0,1] \)
  and all \( x \), then the inverse \( (A(1,\cdot))^{-1} \) has integrable positive log-norm.
\end{lemma}
\begin{proof}\leanok
  \uses{FlowCocycle}
  Identical to \ref{integrableLogNorm_timeOne}, inserting the measurable matrix-inversion map
  and evaluating the dominating hypothesis at \( s = 1 \).
\end{proof}

\begin{proposition}[Discrete filtration for the time-one data]\label{exists_isOseledetsFiltration_timeOne}
  \lean{Oseledets.exists_isOseledetsFiltration_timeOne}\leanok
  Let \( \mu \) be a probability measure, \( \varphi \) a measure-preserving flow with
  \( \varphi(1) \) ergodic, and \( A \) a flow cocycle with the two uniform integrable
  dominators \( g, g' \). Then there exist \( k \), exponents \( \lambda : \mathrm{Fin}\,k \to \R \),
  and a subspace family \( V \) forming an Oseledets filtration for the generator \( A(1,\cdot) \)
  over the dynamics \( \varphi(1) \).
\end{proposition}
\begin{proof}\leanok
  \uses{integrableLogNorm_timeOne,integrableLogNorm_timeOne_inv,oseledets_filtration}
  Apply the discrete theorem \ref{oseledets_filtration} to the ergodic map \( \varphi(1) \) and
  generator \( A(1,\cdot) \), feeding it invertibility (\( \det A(1,\cdot) \ne 0 \)),
  measurability, and the two integrability inputs from \ref{integrableLogNorm_timeOne} and
  \ref{integrableLogNorm_timeOne_inv}. The output is the desired discrete Oseledets datum.
\end{proof}

\section{Between integer times}

The discrete theorem controls growth only along integer times; the next results bridge the gap
to a continuous parameter. Both rest on the orbital sublinearity of an integrable function along
the flow's integer orbit.

\begin{lemma}[Error sublinearity along the integer orbit]\label{ae_tendsto_flowError_zero}
  \lean{Oseledets.ae_tendsto_flowError_zero}\leanok
  For integrable \( g, g' \) and a measure-preserving flow \( \varphi \), for almost every
  \( x \) the combined fluctuation along the integer orbit vanishes:
  \[ n^{-1}\bigl(g(\varphi(n)x) + g'(\varphi(n)x)\bigr) \to 0 \quad (n \to \infty). \]
\end{lemma}
\begin{proof}\leanok
  \uses{natCast_eq_iterate}
  Apply the Birkhoff orbital-tail estimate (a.e.\ \( n^{-1} h(T^{[n]}x) \to 0 \) for integrable
  \( h \)) to \( h = g + g' \) and the time-one map \( T = \varphi(1) \), then rewrite the
  iterate orbit \( (\varphi(1))^{[n]}x \) as \( \varphi(n)x \) using \ref{natCast_eq_iterate}.
\end{proof}

\begin{theorem}[Between-times sandwich: continuous growth equals integer-time growth]\label{tendsto_log_norm_atTop_of_discrete}
  \lean{Oseledets.tendsto_log_norm_atTop_of_discrete}\leanok
  Fix a flow \( \varphi \), a flow cocycle \( A \) with uniform forward/inverse one-step controls
  \( g, g' \) on \( [0,1] \), a point \( x \), and a nonzero vector \( v \). Suppose the integer
  fluctuation error vanishes and the integer-time average converges,
  \( n^{-1}\log\norm{A(n,x)v} \to L \). Then the continuous-time average converges to the same
  limit:
  \[ t^{-1}\log\norm{A(t,x)v} \to L \quad (t \to \infty). \]
\end{theorem}
\begin{proof}\leanok
  \uses{FlowCocycle}
  Write \( t = r + n \) with \( n = \lfloor t \rfloor \ge 1 \) and \( r \in [0,1) \). The cocycle
  identity splits \( A(t,x) = A(r, \varphi(n)x)\, A(n,x) \), so with \( w = A(n,x)v \) one
  sandwiches \( \log\norm{A(t,x)v} \) between
  \( \log\norm{w} - \log\norm{(A(r,\varphi(n)x))^{-1}} \) and
  \( \log\norm{w} + \log\norm{A(r,\varphi(n)x)} \). Both correction terms are bounded in absolute
  value by \( g(\varphi(n)x) + g'(\varphi(n)x) \) via the \( [0,1] \) controls (using
  \( \norm M \norm{M^{-1}} \ge 1 \)). Dividing by \( t \), the error term vanishes by hypothesis,
  the discrete average converges to \( L \) along \( \lfloor t \rfloor \), and the floor ratio
  \( \lfloor t \rfloor / t \to 1 \); a squeeze over the floor delivers the continuous-time limit.
\end{proof}

\section{Equivariance at every real time}

The discrete theorem gives equivariance one integer step at a time. To obtain equivariance at
every real time \( t_0 \) we use the intrinsic growth characterization: a vector lies in level
\( V_i \) iff it is zero or its discrete growth \( \limsup \) is \( \le \lambda_i \). The fixed
matrix \( A(t_0,x) \) is a bounded bijection, so it perturbs the per-step log-norm by \( o(n) \),
hence leaves the growth \( \limsup \) unchanged.

\begin{lemma}[Fixed-time log-norm is sublinear]\label{ae_tendsto_logNorm_fixedTime_zero}
  \lean{Oseledets.ae_tendsto_logNorm_fixedTime_zero}\leanok
  Fix a real time \( t_0 \). For almost every \( x \), both
  \( n^{-1}\log\norm{A(t_0, \varphi(n)x)} \) and \( n^{-1}\log\norm{(A(t_0, \varphi(n)x))^{-1}} \)
  tend to \( 0 \) as \( n \to \infty \).
\end{lemma}
\begin{proof}\leanok
  \uses{FlowCocycle,natCast_eq_iterate}
  One builds an integrable function \( H \) dominating both \( \log^{+}\norm{A(t_0,\cdot)} \) and
  \( \log^{+}\norm{(A(t_0,\cdot))^{-1}} \): by induction on \( \lfloor t_0 \rfloor \) one splits
  \( A((\rho+n)+1,\cdot) = A(\rho+n,\varphi(1)\cdot)\,A(1,\cdot) \), bounds \( \log^{+} \) of a
  product by the sum, and uses measure-preservation of \( \varphi(1) \); negative times reduce to
  positive ones via \( A(t_0,y) = (A(-t_0,\varphi(t_0)y))^{-1} \). The two-sided bound
  \( \abs{\log\norm M} \le \log^{+}\norm M + \log^{+}\norm{M^{-1}} \) turns \( H \) into an
  integrable dominator for the absolute log-norm; the Birkhoff orbital tail of \( H \) along the
  integer orbit and a squeeze finish the proof.
\end{proof}

\begin{theorem}[Shift-invariance of the growth limsup]\label{glim_shift}
  \lean{Oseledets.glim_shift}\leanok
  Fix a real time \( t_0 \). For almost every \( x \) and every test vector \( u \), the
  discrete-time growth \( \limsup \) of the cocycle applied to \( u \) at \( x \) equals that at
  \( \varphi(t_0)x \) applied to the pushed-forward vector \( A(t_0,x)\, u \).
\end{theorem}
\begin{proof}\leanok
  \uses{ae_tendsto_logNorm_fixedTime_zero,toCocycle_eq}
  First, a.e.\ the discrete growth average \( n^{-1}\log\norm{\mathrm{cocycle}\,n\,x\,u} \) has
  bounded range, with upper and lower bounds from the Furstenberg--Kesten Fekete inequalities
  \( \log\norm{\mathrm{cocycle}} \le \mathrm{birkhoffSum}(\log^{+}\norm{A(1,\cdot)}) \) (and the
  inverse version), whose Birkhoff averages converge by the ergodic theorem. The cocycle identity
  and \ref{toCocycle_eq} give the shift relation
  \( \mathrm{cocycle}\,n\,(\varphi(t_0)x)\cdot A(t_0,x) = A(t_0,\varphi(n)x)\cdot \mathrm{cocycle}\,n\,x \),
  so the two growth averages differ by \( n^{-1}\bigl(\log\norm{A(t_0,\varphi(n)x)(\cdots)} - \log\norm{\cdots}\bigr) \),
  which is squeezed to \( 0 \) by \ref{ae_tendsto_logNorm_fixedTime_zero}. A difference tending to
  \( 0 \) between two range-bounded sequences leaves the \( \limsup \) unchanged.
\end{proof}

\begin{theorem}[Flow-equivariance of the filtration at every real time]\label{ae_flow_equivariant}
  \lean{Oseledets.ae_flow_equivariant}\leanok
  Let \( V \) be the Oseledets filtration of the time-one data, and fix \( t_0 \in \R \). Then for
  almost every \( x \) and every level \( i \),
  \[ \mathrm{map}\,(A(t_0,x))\,(V_i\, x) = V_i\,(\varphi(t_0)x). \]
\end{theorem}
\begin{proof}\leanok
  \uses{exists_isOseledetsFiltration_timeOne,glim_shift}
  Use the growth characterization \( v \in V_i\,x \iff v = 0 \lor \limsup \le \lambda_i \) at
  \( x \) and, pulled back along the measure-preserving \( \varphi(t_0) \), at \( \varphi(t_0)x \).
  The map \( P = A(t_0,x) \) is a bijection with inverse \( A(t_0,x)^{-1} \). For a non-bottom
  level, prove both inclusions by \( \mathrm{le\_antisymm} \): membership of \( v \) (resp.\
  \( P^{-1}v \)) is transported through the equivalence of growth \( \limsup \)s supplied by
  \ref{glim_shift}. The bottom level \( V_k = \bot \) maps to \( \bot \).
\end{proof}

\section{The continuous-flow theorem}

\begin{theorem}[Continuous-flow multiplicative ergodic theorem]\label{oseledets_flow}
  \lean{Oseledets.oseledets_flow}\leanok
  Let \( \mu \) be a probability measure on \( X \), let \( \varphi \) be a measure-preserving
  one-parameter flow whose time-one map \( \varphi(1) \) is \( \mu \)-ergodic, and let \( A \) be a
  continuous-time linear cocycle over \( \varphi \) valued in invertible \( d \times d \) real
  matrices. Suppose \( g, g' \in L^1(\mu) \) satisfy, for all \( s \in [0,1] \) and all \( x \),
  \[ \log^{+}\norm{A(s,x)} \le g(x), \qquad \log^{+}\norm{(A(s,x))^{-1}} \le g'(x). \]
  Then there exist \( k \in \N \), a strictly decreasing sequence of exponents
  \( \lambda : \mathrm{Fin}\,k \to \R \), and a measurable family of subspaces
  \( V : \mathrm{Fin}(k+1) \to X \to \mathrm{Submodule}\,\R\,(\mathrm{EuclideanSpace}\,\R\,(\mathrm{Fin}\,d)) \)
  such that:
  \begin{itemize}
    \item \( \lambda \) is strictly decreasing and each \( V_i \) is a measurable subspace family;
    \item \textbf{(full flow equivariance)} for every \( t \in \R \), almost every \( x \) has
      \( A(t,x) \cdot V_i\,x = V_i\,(\varphi(t)x) \) for all \( i \);
    \item almost every \( x \) carries the strict flag
      \( \top = V_0\,x \supsetneq \cdots \supsetneq V_k\,x = \bot \), and on each stratum
      \( v \in V_i\,x \setminus V_{i+1}\,x \) the continuous-time growth rate is exactly
      \( \lambda_i \):
      \[ t^{-1}\log\norm{A(t,x)v} \to \lambda_i \quad (t \to \infty). \]
  \end{itemize}
\end{theorem}
\begin{proof}\leanok
  \uses{exists_isOseledetsFiltration_timeOne,ae_flow_equivariant,ae_tendsto_flowError_zero,tendsto_log_norm_atTop_of_discrete,toCocycle_eq}
  Take \( (k, \lambda, V) \) from the reduction \ref{exists_isOseledetsFiltration_timeOne}; this
  supplies strict antitonicity, measurability, the strict flag, and the integer-time growth rates.
  Full flow equivariance at each \( t \) is \ref{ae_flow_equivariant}. For the continuous-time
  growth, work a.e.\ on the discrete-conclusion set intersected with the error-sublinearity set of
  \ref{ae_tendsto_flowError_zero}. Fix a stratum vector \( v \neq 0 \); rewriting via the reduction
  identity \ref{toCocycle_eq} turns the discrete stratum growth
  \( n^{-1}\log\norm{\mathrm{cocycle}\,n\,x\,v} \to \lambda_i \) into
  \( n^{-1}\log\norm{A(n,x)v} \to \lambda_i \). The between-times sandwich
  \ref{tendsto_log_norm_atTop_of_discrete} then upgrades this integer-time limit to the
  continuous-parameter limit \( t^{-1}\log\norm{A(t,x)v} \to \lambda_i \).
\end{proof}
